% ============================================================
% Section 5: Results
% ============================================================
\section{Results}

% [PLACEHOLDER: Values marked [X] require experimental data.
%  All tables structured for expected results format.]

\subsection{Main Results}

\subsubsection{Dental Fluorosis (DF)}

\begin{table}[htbp]
\centering
\caption{Main results on DF test set (n=60). Best in \textbf{bold}, second \underline{underlined}. $\dagger$: paper-reported.}
\label{tab:df_main}
\begin{tabular}{lccccccc}
\toprule
\textbf{Method} & \textbf{Acc $\uparrow$} & \textbf{F1 $\uparrow$} & \textbf{QWK $\uparrow$} & \textbf{ECE $\downarrow$} & \textbf{\%Uni $\uparrow$} & \textbf{U-ECE $\downarrow$} & \textbf{AUROC($u$) $\uparrow$} \\
\midrule
ViT-B + CE & [X] & [X] & [X] & [X] & [X] & [X] & -- \\
MLTrMR$^\dagger$ \cite{xu2025mltrmr} & 80.19 & 75.79 & -- & -- & -- & -- & -- \\
\midrule
ViT-B + EDL Only & [X] & [X] & [X] & [X] & [X] & [X] & [X] \\
ViT-B + ORCU Only & [X] & [X] & [X] & [X] & [X] & [X] & [X] \\
ViT-B + EDL+ORCU (Ours) & [X] & [X] & [X] & [X] & [X] & [X] & [X] \\
\bottomrule
\end{tabular}
\end{table}

\subsubsection{Skeletal Fluorosis (SF)}

\begin{table}[htbp]
\centering
\caption{Main results on SF test set (n=24).}
\label{tab:sf_main}
\begin{tabular}{lccccccc}
\toprule
\textbf{Method} & \textbf{Acc $\uparrow$} & \textbf{F1 $\uparrow$} & \textbf{QWK $\uparrow$} & \textbf{ECE $\downarrow$} & \textbf{\%Uni $\uparrow$} & \textbf{U-ECE $\downarrow$} & \textbf{AUROC($u$) $\uparrow$} \\
\midrule
ResNet-50 + CE & [X] & [X] & [X] & [X] & [X] & [X] & -- \\
Mwinc-Mamba$^\dagger$ \cite{xu2025mwinc} & 66.67 & -- & -- & -- & -- & -- & -- \\
5 Radiologists Avg$^\dagger$ & 70.00 & -- & -- & -- & -- & -- & -- \\
\midrule
ResNet-50 + EDL Only & [X] & [X] & [X] & [X] & [X] & [X] & [X] \\
ResNet-50 + ORCU Only & [X] & [X] & [X] & [X] & [X] & [X] & [X] \\
ResNet-50 + EDL+ORCU (Ours) & [X] & [X] & [X] & [X] & [X] & [X] & [X] \\
ResNet-50 + EDL+ORCU + MR (Ours) & [X] & [X] & [X] & [X] & [X] & [X] & [X] \\
\bottomrule
\end{tabular}
\end{table}

\subsection{Ablation Results}

\begin{table}[htbp]
\centering
\caption{Ablation study results.}
\label{tab:ablation}
\begin{tabular}{lcccc}
\toprule
\textbf{Variant} & \textbf{Acc} & \textbf{QWK} & \textbf{ECE} & \textbf{\%Uni} \\
\midrule
\multicolumn{5}{c}{\textit{DF (ViT-B)}} \\
CE (baseline) & [X] & [X] & [X] & [X] \\
EDL Only & [X] & [X] & [X] & [X] \\
ORCU Only & [X] & [X] & [X] & [X] \\
EDL+ORCU ($\lambda=0.5$) & [X] & [X] & [X] & [X] \\
\midrule
\multicolumn{5}{c}{\textit{A3: $\lambda$ sweep (DF)}} \\
$\lambda=0.1$ & [X] & [X] & [X] & [X] \\
$\lambda=0.3$ & [X] & [X] & [X] & [X] \\
$\lambda=0.5$ (default) & [X] & [X] & [X] & [X] \\
$\lambda=0.7$ & [X] & [X] & [X] & [X] \\
$\lambda=1.0$ & [X] & [X] & [X] & [X] \\
\midrule
\multicolumn{5}{c}{\textit{A6: Label type (SF)}} \\
One-hot majority vote & [X] & [X] & [X] & [X] \\
Multi-rater soft label & [X] & [X] & [X] & [X] \\
\bottomrule
\end{tabular}
\end{table}

\subsection{Uncertainty Analysis}

\textbf{Evidence per class (U1):} [Expected: Moderate/Severe classes show lower evidence due to fewer samples, especially SF.]

\textbf{Uncertainty vs error (U2):} AUROC of $u$ for misclassification: DF=[X], SF=[X].

\textbf{OOD detection (U3):} [Expected: EDL assigns significantly higher $u$ to OOD samples.]

\textbf{Uncertainty by body part (U4, SF):} [Expected: Pelvis views show highest $u$.]

\textbf{Rater agreement vs uncertainty (U5):} Spearman correlation between inter-rater entropy and model $u$: [X] ($p$=[X]).

\textbf{RecAcc \& SafePred (U6):} At $\theta=0.30$, [X]\% flagged for review (RecAcc=[X]\%), [X]\% safe (SafePred=[X]\%).

\subsection{Clinical Report Generation}

Example auto-generated diagnosis reports for representative DF/SF cases are shown in the supplementary material. The reports, produced from EDL outputs without additional training, include predicted grade, confidence, uncertainty level, per-class evidence, and clinical recommendations stratified by uncertainty.
